\section{Results}
\label{sec:results}

\subsection{Typed-program supervision more than doubles strict accuracy}

Table \ref{tab:pacs-main} answers RQ1.  The untuned Qwen3-8B model solves
\PacsBaseCorrect/\PacsN cases (\PacsBaseA\%).  The cloud planner improves this to
\PacsTeacherCorrect/\PacsN (\PacsTeacherA\%).  The typed-program model solves
\PacsOursCorrectA/\PacsN (\PacsOursA\%), an absolute gain of \PacsGainBase{} percentage points
over the same base model and \PacsGainTeacher{} points over the cloud planner.  The Wilson
intervals do not overlap either baseline interval.

\begin{table}[t]
  \caption{PACS strict execution accuracy.  All Channel-A rows use the same \PacsN{} cases and a
  single temperature-zero free decode.  ``Valid tree'' is reported for diagnosis only: explicit
  abstentions do not emit a tree, and validity does not imply correctness.}
  \label{tab:pacs-main}
  \centering
  \small
  \begin{tabular}{llrrrcc}
    \toprule
    System & Surface & Correct & $n$ & Strict acc. & 95\% Wilson CI & Valid tree \\
    \midrule
    Untuned Qwen3-8B & A & 332 & 922 & 36.01 & [32.97, 39.16] & 77.77 \\
    Cloud planner & A & 450 & 922 & 48.81 & [45.59, 52.03] & 52.28 \\
    \textbf{Typed-program Qwen3-8B} & A & \textbf{687} & 922 & \textbf{74.51}
      & \textbf{[71.60, 77.22]} & 71.58 \\
    Typed-program Qwen3-8B & B & 675 & 922 & 73.21 & [70.26, 75.97] & 70.50 \\
    \bottomrule
  \end{tabular}
\end{table}

Tree-valid rate alone ranks the untuned base above the specialised model, despite a 38.50-point
gap in strict accuracy.  This is direct evidence for the distinction in Table
\ref{tab:guarantee-boundary}: a syntactically and type-valid program may still reverse roles,
select the wrong field, or express the wrong scope.  Conversely, an explicit abstention is a
scoreable final action without a tree.

\subsection{Surface naturalisation causes a small aggregate drop}

For RQ2, the specialised model falls from \PacsOursA\% on Channel A to \PacsOursB\% on Channel
B, a difference of \PacsSurfaceDrop{} percentage points.  Macro-family accuracy falls from
\PacsMacroA\% to \PacsMacroB\%.  This suggests that the parser is not merely reproducing one
fixed wording, but the result should not be interpreted as unrestricted linguistic robustness:
both surfaces preserve the same literals and are generated under explicit logic gates.

The largest A--B family decline is 3.38 points on relational composition (F6); supplier
concentration (F3) instead improves by 2.96 points.  A paired item-level significance test is
left out of the present draft because the intended strict Boolean-matching fix must first be
applied uniformly to every scorer and the paired decisions regenerated from the frozen outputs.

\subsection{Improvements are concentrated in compositional set operations}

Table \ref{tab:pacs-family} answers RQ3.  The specialised parser is strongest on overlap and
exclusion (F5, 94.49\%) and improves the base by more than 60 points.  It also substantially
improves cross-buyer comparison (F4) and relational composition (F6).  Supplier concentration
(F3) and disclosure/compliance (F7) remain near 50\% and 42\%, respectively; on both families the
cloud model is competitive with or slightly above the specialised model.  These failures are
obscured by the micro average, particularly because F7 contains only 67 cases.

\begin{table}[t]
  \caption{Strict accuracy by PACS family.  Exact correct counts are given in parentheses.
  Surface B is available only for the specialised parser.}
  \label{tab:pacs-family}
  \centering
  \small
  \begin{tabular}{llrrrr}
    \toprule
    & Family ($n$) & Untuned A & Cloud A & Typed A & Typed B \\
    \midrule
    F1 & Spending/volume (148) & 41.89 (62) & 62.84 (93) & \textbf{84.46 (125)} & 81.08 (120) \\
    F2 & Temporal change (149) & 57.72 (86) & 72.48 (108) & \textbf{83.22 (124)} & 81.88 (122) \\
    F3 & Supplier concentration (135) & 19.26 (26) & 51.85 (70) & 51.11 (69) & \textbf{54.07 (73)} \\
    F4 & Cross-buyer comparison (147) & 39.46 (58) & 44.22 (65) & \textbf{81.63 (120)} & 79.59 (117) \\
    F5 & Overlap/exclusion (127) & 32.28 (41) & 33.07 (42) & \textbf{94.49 (120)} & 93.70 (119) \\
    F6 & Relational composition (149) & 28.19 (42) & 28.86 (43) & \textbf{67.79 (101)} & 65.10 (97) \\
    F7 & Disclosure/compliance (67) & 25.37 (17) & \textbf{43.28 (29)} & 41.79 (28) & 40.30 (27) \\
    \midrule
    & Macro family & 34.88 & 48.09 & \textbf{72.07} & 70.82 \\
    \bottomrule
  \end{tabular}
\end{table}

Depth diagnostics are not monotonic.  The typed parser obtains 92.11\% at L1, 61.06\% at L2,
and 73.36\% at L3.  The lower L2 value is caused by the distribution of families and status
types rather than proving that two-step programs are intrinsically harder than L3.  Declared
seen and unseen shapes score 74.01\% and 74.95\%, respectively.  This near equality is
encouraging, but it is not by itself proof of CFQ-style compositional generalisation because the
split lacks an atom/compound-divergence construction and its near-duplicate screen is
subsampled.

\subsection{The full runtime improves on the balanced system snapshot}

Table \ref{tab:runtime-main} answers RQ4 on the separate \RuntimeN{}-item snapshot.  The full local
runtime solves \RuntimeOursCorrect/\RuntimeN cases (\RuntimeOurs\%), compared with
\RuntimeBaseCorrect/\RuntimeN (\RuntimeBase\%) for the untuned local system and
\RuntimeTeacherCorrect/\RuntimeN (\RuntimeTeacher\%) for the stored cloud-teacher run.  The
largest category differences relative to the untuned system occur on bridge joins (14/20 versus
2/20), comparisons (17/20 versus 7/20), and categorical questions (13/20 versus 8/20).  All three
systems score 20/20 on the snapshot's sum category; because the underlying data contain
unnormalised non-GBP additive rows, we do not interpret that category as evidence of
currency-correct aggregation.

\begin{table}[t]
  \caption{Secondary full-runtime result on a balanced frozen artifact (13 categories,
  20 cases each).  The incomplete external run manifest makes this weaker reproducibility
  evidence than Table \ref{tab:pacs-main}.}
  \label{tab:runtime-main}
  \centering
  \small
  \begin{tabular}{lrrrc}
    \toprule
    System & Correct & $n$ & Accuracy & 95\% Wilson CI \\
    \midrule
    Untuned local runtime & 183 & 260 & 70.38 & [64.57, 75.60] \\
    Cloud-teacher runtime & 192 & 260 & 73.85 & [68.18, 78.81] \\
    \textbf{Full local runtime} & \textbf{224} & 260 & \textbf{86.15} & \textbf{[81.43, 89.83]} \\
    \bottomrule
  \end{tabular}
\end{table}

\subsection{Qualitative bridge case}

One PACS F6--L3 case asks for the buyers reached by first finding the lead suppliers of an anchor
buyer, following those supplier values to other notices, and excluding the exact anchor string.
The untuned model filters \texttt{supplier\_name} by the buyer literal and projects
\texttt{contract\_node\_id}; its well-typed set difference executes to the empty set.  The
specialised parser instead emits Equation \ref{eq:bridge-example} and returns the 17 stored
buyer-name strings in the oracle.  The output includes surface variants such as
``EDUCATION AUTHORITY'', ``Education Authority NI'', and ``the Education Authority'', which is
useful evidence of both the compositional success and the unresolved entity-name limitation.

This case belongs to the direct typed-program protocol.  It demonstrates semantic parsing and
deterministic set execution, but it is not evidence that the deployment understander,
conditional planner, or release checks ran for that question.  A Figure-1 end-to-end example
must therefore be exported from one full runtime trace rather than reusing this PACS row.
